\section{Conclusiones}

El análisis numérico desarrollado demuestra rigurosamente la convergencia entre la dinámica estocástica microscópica y los modelos de transporte macroscópicos.

El estudio de caminos aleatorios en 1D, 2D y 3D valida empíricamente la ley de difusión de Fick. La evolución temporal del desplazamiento cuadrático medio confirma la relación lineal $\langle r^2 \rangle \propto t$, demostrando ser invariante ante cambios de dimensionalidad o geometría de coordenadas. Las desviaciones fraccionales observadas ($\approx 2-3\%$) frente a la pendiente teórica ideal constituyen fluctuaciones estadísticas esperables dada la finitud del ensamble ($N=1000$).

Al evaluar la dinámica espacial, la correspondencia entre la discretización estocástica y la ecuación determinista de difusión prueba que el límite macroscópico de los caminantes aleatorios reproduce la dinámica de ensanchamiento gaussiano continuo. El análisis de la evolución entrópica del campo escalar de densidades constata una rápida aproximación hacia el estado asintótico de máximo mezclado, quedando este proceso de termalización gobernado por la relación de escalamiento $t_{\mathrm{eq}} \propto L^2$ para el tiempo de relajación.

Por último, el modelo de percolación evidencia los límites de validez de esta difusión isotrópica regular. Al alcanzar la densidad crítica $p_c$, el sistema experimenta una transición de fase de segundo orden; la geometría fractal del cúmulo invasor ($D_f < 2$) introduce heterogeneidades espaciales que colapsan el transporte clásico, forzando un régimen de difusión anómala característico de los medios porosos.